\documentclass[a4paper,10pt]{article}
\usepackage[top=14mm,bottom=16mm,left=15mm,right=15mm]{geometry}
\usepackage{fontspec}
\usepackage{xeCJK}
\setmainfont{Times New Roman}
\setCJKmainfont{Yu Gothic}
\usepackage{array,tabularx}
\setlength{\parindent}{0pt}
\newcommand{\answerrow}[2]{%
\textbf{#1}\quad #2\\[-1mm]
\fbox{\parbox[t][34mm][t]{0.96\linewidth}{\hspace{1mm}}}\par\vspace{3mm}}
\begin{document}
\begin{center}
{\LARGE\bfseries ソーラーカーMPC-EMS 完全処理フロー体験書}\\[2mm]
{\Large 解答用紙}\\[5mm]
\begin{tabularx}{\linewidth}{|p{25mm}|X|p{28mm}|}
\hline
氏名 & & 採点\\ \hline
実施日 & & 確認者\\ \hline
\end{tabularx}
\end{center}
\vspace{4mm}
\answerrow{Q1}{夜間補機}
\answerrow{Q2}{起動}
\answerrow{Q3}{UTC}
\answerrow{Q4}{filter}
\answerrow{Q5}{距離・勾配}
\answerrow{Q6}{nominal/risk}
\answerrow{Q7}{補機の電気側加算}
\answerrow{Q8}{pack・SoC}
\answerrow{Q9}{夜間pack}
\answerrow{Q10}{CEM / L-BFGS-B}
\answerrow{Q11}{replan}
\answerrow{Q12}{rate limit}
\answerrow{Q13}{再現性}
\answerrow{Q14}{timeout}
\answerrow{Q15}{MLE検証}
\answerrow{Q16}{全周期}
\answerrow{Q17}{満充電張り付き}
\answerrow{Q18}{3026.9 km}
\answerrow{Q19}{配布同期}
\answerrow{Q20}{release gate}
\answerrow{Q21}{control展開・区間時間}
\answerrow{Q22}{空力・転がり・勾配・車輪出力}
\answerrow{Q23}{drive map二線形補間}
\answerrow{Q24}{DC・pack・IV完全計算}
\answerrow{Q25}{6段SoC表・終端使用可能Wh}
\answerrow{Q26}{終端SoC許容帯}
\answerrow{Q27}{有限差分勾配・一反復}
\answerrow{Q28}{L-BFGS二ループ}
\answerrow{Q29}{有限格子全列挙証明}
\answerrow{Q30}{CEMと数学的証明の違い}
\answerrow{Q31}{detail CSV収支不整合の調査}
\answerrow{Q32}{日射APIの時刻定義}
\answerrow{Q33}{時刻・距離の双線形補間}
\answerrow{Q34}{70 km/h車両総負荷}
\answerrow{Q35}{冬季PVと電池電力}
\answerrow{Q36}{外側端数と1 Hz detail}
\answerrow{Q37}{回生効率と回生利用率}
\answerrow{Q38}{同期RMSEと区間energy}
\end{document}
